\section{账本事件的制度与社会深描：tang}

\label{sec:ledger-depth-tang}

本组以账本所列事件为可核锚点，转而追问军事、财政、地方行政、宗教与跨区域网络怎样在具体空间中相互限制。官方叙述、文书和物质遗存各有覆盖范围，不能由任何一种来源替代普通家庭与流动人群的完整经验。

\subsection{区域制度与资源}

条目\texttt{ST-581-SUI-ACCESSION}所记的开皇元年（581）二月“杨坚受北周静帝禅，建国号隋，改元开皇”，发生在隋的政治语境中。账本将其原因概括为杨坚掌握禁军与中枢，宇文氏政权继承基础已被削弱，并提示隋整合关中军政资源，开始再统一进程。这类节点进入地方社会时，要经过官员、军队、书吏、商人、寺院、道路与家庭等不同中介；同一政策或战事在城市、边地和乡村不会具有相同的节奏。

核心来源为《隋书·高祖纪上》；《北史·隋本纪上》；《资治通鉴》隋纪（开皇元年二月条）。它们可以较有把握地支持年序、制度公布、使节活动或特定地点的记载，却不能自动给出人口、财富、兵额、族属和内心动机的完整调查。账本的证据注记是“纪年可靠；《隋书》为唐初官修，受禅仪文与动机含胜者叙事”。因此，不能将官府的分类、战报的数字、宗教文本的愿望或后出的总括直接写成全体社会事实。

更合适的解释是，将这一事件看成资源、信息和权力重新连接的一环：命令是否抵达，取决于交通和地方协作；征发是否落实，取决于登记、市场和劳力；文化和宗教网络能否延续，也受战争、迁徙和赞助条件影响。现有材料只照亮其中一部分，叙事必须让区域差异与材料沉默继续可见。

条目\texttt{ST-582-TURKIC-DEFENSE}所记的开皇二年（582）“沙钵略可汗联结北周旧部南侵，隋以边防与分化策略应对”，发生在隋与突厥的政治语境中。账本将其原因概括为突厥利用隋周易代的边防空隙，试图维持北方军事压力，并提示隋将北方威胁与南征分开处理，分化突厥成为可用外交工具。这类节点进入地方社会时，要经过官员、军队、书吏、商人、寺院、道路与家庭等不同中介；同一政策或战事在城市、边地和乡村不会具有相同的节奏。

核心来源为《隋书·高祖纪上》；《隋书·突厥传》；《北史·突厥传》。它们可以较有把握地支持年序、制度公布、使节活动或特定地点的记载，却不能自动给出人口、财富、兵额、族属和内心动机的完整调查。账本的证据注记是“年序可靠；突厥内部联系主要见汉文朝廷叙事，须警惕归顺话语”。因此，不能将官府的分类、战报的数字、宗教文本的愿望或后出的总括直接写成全体社会事实。

更合适的解释是，将这一事件看成资源、信息和权力重新连接的一环：命令是否抵达，取决于交通和地方协作；征发是否落实，取决于登记、市场和劳力；文化和宗教网络能否延续，也受战争、迁徙和赞助条件影响。现有材料只照亮其中一部分，叙事必须让区域差异与材料沉默继续可见。

条目\texttt{ST-583-KAIHUANG-CODE}所记的开皇三年（583）“隋廷颁行修定后的开皇律令，减省死刑条目并重整刑名”，发生在隋的政治语境中。账本将其原因概括为北周旧律繁重，新朝需要以成文法统合关东、关中治理传统，并提示开皇律成为唐律的重要前史，实际审判仍须以文书和个案检验。这类节点进入地方社会时，要经过官员、军队、书吏、商人、寺院、道路与家庭等不同中介；同一政策或战事在城市、边地和乡村不会具有相同的节奏。

核心来源为《隋书·刑法志》；《北史·刑法志》；《通典·刑法》。它们可以较有把握地支持年序、制度公布、使节活动或特定地点的记载，却不能自动给出人口、财富、兵额、族属和内心动机的完整调查。账本的证据注记是“制度颁布可靠；后出志书所录条文不等于全国地方执行”。因此，不能将官府的分类、战报的数字、宗教文本的愿望或后出的总括直接写成全体社会事实。

更合适的解释是，将这一事件看成资源、信息和权力重新连接的一环：命令是否抵达，取决于交通和地方协作；征发是否落实，取决于登记、市场和劳力；文化和宗教网络能否延续，也受战争、迁徙和赞助条件影响。现有材料只照亮其中一部分，叙事必须让区域差异与材料沉默继续可见。

条目\texttt{ST-588-SOUTHERN-CAMPAIGN}所记的开皇八年（588）十月“隋分水陆诸军南下，发动灭陈战争”，发生在隋与陈的政治语境中。账本将其原因概括为陈朝军政财政衰弱，隋已完成北方整合并取得北边缓冲，并提示长江防线被多路进攻穿透，统一战争进入决定阶段。这类节点进入地方社会时，要经过官员、军队、书吏、商人、寺院、道路与家庭等不同中介；同一政策或战事在城市、边地和乡村不会具有相同的节奏。

核心来源为《隋书·高祖纪下》；《隋书·陈后主纪》；《资治通鉴》隋纪（开皇八年十月条）。它们可以较有把握地支持年序、制度公布、使节活动或特定地点的记载，却不能自动给出人口、财富、兵额、族属和内心动机的完整调查。账本的证据注记是“纪年可靠；兵力与行军路线多取自官修战功叙事，数字不宜实数化”。因此，不能将官府的分类、战报的数字、宗教文本的愿望或后出的总括直接写成全体社会事实。

更合适的解释是，将这一事件看成资源、信息和权力重新连接的一环：命令是否抵达，取决于交通和地方协作；征发是否落实，取决于登记、市场和劳力；文化和宗教网络能否延续，也受战争、迁徙和赞助条件影响。现有材料只照亮其中一部分，叙事必须让区域差异与材料沉默继续可见。

条目\texttt{ST-589-CHEN-CONQUEST}所记的开皇九年（589）正月“隋军入建康，陈后主被俘，陈朝覆亡”，发生在隋与陈的政治语境中。账本将其原因概括为隋军渡江后陈廷指挥失灵，江南守军难以形成持久协同，并提示南北分立结束；隋须处理江南户籍、叛乱风险和行政整合。这类节点进入地方社会时，要经过官员、军队、书吏、商人、寺院、道路与家庭等不同中介；同一政策或战事在城市、边地和乡村不会具有相同的节奏。

核心来源为《隋书·高祖纪下》；《隋书·陈后主纪》；《陈书·后主纪》；《资治通鉴》隋纪（开皇九年正月条）。它们可以较有把握地支持年序、制度公布、使节活动或特定地点的记载，却不能自动给出人口、财富、兵额、族属和内心动机的完整调查。账本的证据注记是“核心事实与月份可靠；俘获与宫廷逸事多经亡国叙事修饰”。因此，不能将官府的分类、战报的数字、宗教文本的愿望或后出的总括直接写成全体社会事实。

更合适的解释是，将这一事件看成资源、信息和权力重新连接的一环：命令是否抵达，取决于交通和地方协作；征发是否落实，取决于登记、市场和劳力；文化和宗教网络能否延续，也受战争、迁徙和赞助条件影响。现有材料只照亮其中一部分，叙事必须让区域差异与材料沉默继续可见。

条目\texttt{ST-590-JIANGNAN-REBELLION}所记的开皇十年（590）“江南旧陈地区多起反隋，杨素等率军镇压并重建地方控制”，发生在隋与江南诸州的政治语境中。账本将其原因概括为新朝清查户口、改置州县并接收陈朝军民，地方对赋役与官吏更替反应强烈，并提示统一由建康易手转为军事镇压和行政编户的地方接管。这类节点进入地方社会时，要经过官员、军队、书吏、商人、寺院、道路与家庭等不同中介；同一政策或战事在城市、边地和乡村不会具有相同的节奏。

核心来源为《隋书·高祖纪下》；《隋书·杨素传》；《资治通鉴》隋纪（开皇十年条）。它们可以较有把握地支持年序、制度公布、使节活动或特定地点的记载，却不能自动给出人口、财富、兵额、族属和内心动机的完整调查。账本的证据注记是“起事与镇压可靠；参与者规模、动机和死亡数字主要来自中央军功叙事”。因此，不能将官府的分类、战报的数字、宗教文本的愿望或后出的总括直接写成全体社会事实。

更合适的解释是，将这一事件看成资源、信息和权力重新连接的一环：命令是否抵达，取决于交通和地方协作；征发是否落实，取决于登记、市场和劳力；文化和宗教网络能否延续，也受战争、迁徙和赞助条件影响。现有材料只照亮其中一部分，叙事必须让区域差异与材料沉默继续可见。

条目\texttt{ST-598-GOGURYEO-FIRST-CAMPAIGN}所记的开皇十八年（598）“隋水陆军征高句丽，因疾疫、风浪和补给困难撤回”，发生在隋与高句丽的政治语境中。账本将其原因概括为高句丽扩张与隋统一后的天下秩序发生摩擦，隋试图以远征确立北东方向支配力，并提示远征暴露辽东战争的运输和气候限制，未改变高句丽的独立行动能力。这类节点进入地方社会时，要经过官员、军队、书吏、商人、寺院、道路与家庭等不同中介；同一政策或战事在城市、边地和乡村不会具有相同的节奏。

核心来源为《隋书·高祖纪下》；《隋书·东夷传·高丽》；《资治通鉴》隋纪（开皇十八年条）。它们可以较有把握地支持年序、制度公布、使节活动或特定地点的记载，却不能自动给出人口、财富、兵额、族属和内心动机的完整调查。账本的证据注记是“出兵与撤军可靠；伤亡、船数和战果均为单方记载，须保留高句丽材料缺口”。因此，不能将官府的分类、战报的数字、宗教文本的愿望或后出的总括直接写成全体社会事实。

更合适的解释是，将这一事件看成资源、信息和权力重新连接的一环：命令是否抵达，取决于交通和地方协作；征发是否落实，取决于登记、市场和劳力；文化和宗教网络能否延续，也受战争、迁徙和赞助条件影响。现有材料只照亮其中一部分，叙事必须让区域差异与材料沉默继续可见。

条目\texttt{ST-604-YANGDI-ACCESSION}所记的仁寿四年（604）七月“杨广即帝位，是为隋炀帝”，发生在隋的政治语境中。账本将其原因概括为文帝晚年继承安排反复，杨广已经营东宫及关陇官僚网络，并提示新君得以推进东都、运河和远征等高动员政策。这类节点进入地方社会时，要经过官员、军队、书吏、商人、寺院、道路与家庭等不同中介；同一政策或战事在城市、边地和乡村不会具有相同的节奏。

核心来源为《隋书·炀帝纪上》；《北史·隋本纪下》；《资治通鉴》隋纪（仁寿四年条）。它们可以较有把握地支持年序、制度公布、使节活动或特定地点的记载，却不能自动给出人口、财富、兵额、族属和内心动机的完整调查。账本的证据注记是“即位日期可靠；夺嫡与文帝死因细节主要来自唐修亡国叙事，不能当作定论”。因此，不能将官府的分类、战报的数字、宗教文本的愿望或后出的总括直接写成全体社会事实。

更合适的解释是，将这一事件看成资源、信息和权力重新连接的一环：命令是否抵达，取决于交通和地方协作；征发是否落实，取决于登记、市场和劳力；文化和宗教网络能否延续，也受战争、迁徙和赞助条件影响。现有材料只照亮其中一部分，叙事必须让区域差异与材料沉默继续可见。

条目\texttt{ST-605-EASTERN-CAPITAL}所记的大业元年（605）“隋廷营建东都洛阳并迁置人口、仓储和官署”，发生在隋的政治语境中。账本将其原因概括为关中供给紧张，东都又居运河与关东交通枢纽，炀帝需更直接控制东部资源，并提示洛阳成为东部行政物流中心，也加重徙民、营建和粮运负担。这类节点进入地方社会时，要经过官员、军队、书吏、商人、寺院、道路与家庭等不同中介；同一政策或战事在城市、边地和乡村不会具有相同的节奏。

核心来源为《隋书·炀帝纪上》；《隋书·宇文恺传》；隋唐洛阳城遗址考古报告。它们可以较有把握地支持年序、制度公布、使节活动或特定地点的记载，却不能自动给出人口、财富、兵额、族属和内心动机的完整调查。账本的证据注记是“工程启动可靠；徙民数和宫城规模多为史书夸饰，须以洛阳城址分期校验”。因此，不能将官府的分类、战报的数字、宗教文本的愿望或后出的总括直接写成全体社会事实。

更合适的解释是，将这一事件看成资源、信息和权力重新连接的一环：命令是否抵达，取决于交通和地方协作；征发是否落实，取决于登记、市场和劳力；文化和宗教网络能否延续，也受战争、迁徙和赞助条件影响。现有材料只照亮其中一部分，叙事必须让区域差异与材料沉默继续可见。

条目\texttt{ST-605-TONGJI-CANAL}所记的大业元年（605）“隋开通通济渠，沟通洛阳、淮河与江都方向水路”，发生在隋的政治语境中。账本将其原因概括为东都建设及南北粮食、军队转运需要稳定水路，既有河道被联结扩展，并提示运河提高跨区域调运能力，也使短期征发成为地方不满来源。这类节点进入地方社会时，要经过官员、军队、书吏、商人、寺院、道路与家庭等不同中介；同一政策或战事在城市、边地和乡村不会具有相同的节奏。

核心来源为《隋书·炀帝纪上》；《隋书·食货志》；隋唐大运河通济渠考古调查报告。它们可以较有把握地支持年序、制度公布、使节活动或特定地点的记载，却不能自动给出人口、财富、兵额、族属和内心动机的完整调查。账本的证据注记是“开凿年份可靠；役夫数量和死亡叙事受亡国框架影响，遗存不能证实具体劳力规模”。因此，不能将官府的分类、战报的数字、宗教文本的愿望或后出的总括直接写成全体社会事实。

更合适的解释是，将这一事件看成资源、信息和权力重新连接的一环：命令是否抵达，取决于交通和地方协作；征发是否落实，取决于登记、市场和劳力；文化和宗教网络能否延续，也受战争、迁徙和赞助条件影响。现有材料只照亮其中一部分，叙事必须让区域差异与材料沉默继续可见。

条目\texttt{ST-607-COMMANDERY-REFORM}所记的大业三年（607）“隋废州改郡，调整地方行政名号和层级”，发生在隋的政治语境中。账本将其原因概括为统一后州郡层级繁多，炀帝试图整齐行政名目并强化中央统摄，并提示郡制成为隋末唐初地方治理背景，唐初复州显示名制不能单独解决统治成本。这类节点进入地方社会时，要经过官员、军队、书吏、商人、寺院、道路与家庭等不同中介；同一政策或战事在城市、边地和乡村不会具有相同的节奏。

核心来源为《隋书·炀帝纪上》；《隋书·地理志》；《通典·州郡》。它们可以较有把握地支持年序、制度公布、使节活动或特定地点的记载，却不能自动给出人口、财富、兵额、族属和内心动机的完整调查。账本的证据注记是“诏改可靠；名称变更与地方实际权力、财政边界不必同步”。因此，不能将官府的分类、战报的数字、宗教文本的愿望或后出的总括直接写成全体社会事实。

更合适的解释是，将这一事件看成资源、信息和权力重新连接的一环：命令是否抵达，取决于交通和地方协作；征发是否落实，取决于登记、市场和劳力；文化和宗教网络能否延续，也受战争、迁徙和赞助条件影响。现有材料只照亮其中一部分，叙事必须让区域差异与材料沉默继续可见。

条目\texttt{ST-609-TUYUHUN-CAMPAIGN}所记的大业五年（609）“隋炀帝亲至河西，隋军击吐谷浑并设置西海、河源等郡”，发生在隋与吐谷浑的政治语境中。账本将其原因概括为吐谷浑控制河西—青海通路并受隋西域政策挤压，炀帝以亲征显示边疆支配力，并提示隋一度把行政名号推入青海高原，高成本驻防预示边防脆弱。这类节点进入地方社会时，要经过官员、军队、书吏、商人、寺院、道路与家庭等不同中介；同一政策或战事在城市、边地和乡村不会具有相同的节奏。

核心来源为《隋书·炀帝纪上》；《隋书·吐谷浑传》；《隋书·地理志》。它们可以较有把握地支持年序、制度公布、使节活动或特定地点的记载，却不能自动给出人口、财富、兵额、族属和内心动机的完整调查。账本的证据注记是“军事行动和置郡见史书；郡治持续时间和实际控制范围须以遗址、文书另核”。因此，不能将官府的分类、战报的数字、宗教文本的愿望或后出的总括直接写成全体社会事实。

更合适的解释是，将这一事件看成资源、信息和权力重新连接的一环：命令是否抵达，取决于交通和地方协作；征发是否落实，取决于登记、市场和劳力；文化和宗教网络能否延续，也受战争、迁徙和赞助条件影响。现有材料只照亮其中一部分，叙事必须让区域差异与材料沉默继续可见。

\subsection{边疆、道路与迁徙}

条目\texttt{ST-611-WANG-BO-REBELLION}所记的大业七年（611）“王薄在长白山起事，隋末山东民变由逃亡转为持续武装活动”，发生在隋与山东起事者的政治语境中。账本将其原因概括为征辽前兵役、粮运和治安压力集中于山东、河北，逃亡者获得山泽据点，并提示中央对东部征发的控制受侵蚀，为瓦岗等势力扩张创造环境。这类节点进入地方社会时，要经过官员、军队、书吏、商人、寺院、道路与家庭等不同中介；同一政策或战事在城市、边地和乡村不会具有相同的节奏。

核心来源为《隋书·炀帝纪下》；《隋书·王薄传》；《资治通鉴》隋纪（大业七年条）。它们可以较有把握地支持年序、制度公布、使节活动或特定地点的记载，却不能自动给出人口、财富、兵额、族属和内心动机的完整调查。账本的证据注记是“王薄起事可确证；民歌与参与者叙事受后世文本保存影响”。因此，不能将官府的分类、战报的数字、宗教文本的愿望或后出的总括直接写成全体社会事实。

更合适的解释是，将这一事件看成资源、信息和权力重新连接的一环：命令是否抵达，取决于交通和地方协作；征发是否落实，取决于登记、市场和劳力；文化和宗教网络能否延续，也受战争、迁徙和赞助条件影响。现有材料只照亮其中一部分，叙事必须让区域差异与材料沉默继续可见。

条目\texttt{ST-612-GOGURYEO-INVASION}所记的大业八年（612）“隋第一次大规模征高句丽，辽东城久攻不下，九军渡辽水后覆败”，发生在隋与高句丽的政治语境中。账本将其原因概括为隋以全国兵役和漕运追求速胜，高句丽依城防、河流和纵深抵消兵力优势，并提示远征耗损军队和财政信誉，山东河北的征发与逃亡压力急剧上升。这类节点进入地方社会时，要经过官员、军队、书吏、商人、寺院、道路与家庭等不同中介；同一政策或战事在城市、边地和乡村不会具有相同的节奏。

核心来源为《隋书·炀帝纪下》；《隋书·东夷传·高丽》；《资治通鉴》隋纪（大业八年条）。它们可以较有把握地支持年序、制度公布、使节活动或特定地点的记载，却不能自动给出人口、财富、兵额、族属和内心动机的完整调查。账本的证据注记是“出征与失败可靠；兵数、伤亡及渡海细节为单方战史，不能视作实数”。因此，不能将官府的分类、战报的数字、宗教文本的愿望或后出的总括直接写成全体社会事实。

更合适的解释是，将这一事件看成资源、信息和权力重新连接的一环：命令是否抵达，取决于交通和地方协作；征发是否落实，取决于登记、市场和劳力；文化和宗教网络能否延续，也受战争、迁徙和赞助条件影响。现有材料只照亮其中一部分，叙事必须让区域差异与材料沉默继续可见。

条目\texttt{ST-613-YANG-XUANGAN}所记的大业九年（613）六月“杨玄感趁炀帝再征高句丽时在黎阳起兵，旋被平定”，发生在隋的政治语境中。账本将其原因概括为征辽使精兵东出，杨玄感又掌黎阳转运，利用京师与前线间的权力空隙，并提示朝廷对重臣、仓储和地方武装的信任瓦解，征辽被迫中断。这类节点进入地方社会时，要经过官员、军队、书吏、商人、寺院、道路与家庭等不同中介；同一政策或战事在城市、边地和乡村不会具有相同的节奏。

核心来源为《隋书·炀帝纪下》；《隋书·杨玄感传》；《资治通鉴》隋纪（大业九年六月条）。它们可以较有把握地支持年序、制度公布、使节活动或特定地点的记载，却不能自动给出人口、财富、兵额、族属和内心动机的完整调查。账本的证据注记是“起兵和覆灭日期可靠；政治口号与同谋范围在纪传间有异说”。因此，不能将官府的分类、战报的数字、宗教文本的愿望或后出的总括直接写成全体社会事实。

更合适的解释是，将这一事件看成资源、信息和权力重新连接的一环：命令是否抵达，取决于交通和地方协作；征发是否落实，取决于登记、市场和劳力；文化和宗教网络能否延续，也受战争、迁徙和赞助条件影响。现有材料只照亮其中一部分，叙事必须让区域差异与材料沉默继续可见。

条目\texttt{ST-614-GOGURYEO-THIRD}所记的大业十年（614）“隋第三次征高句丽，高句丽遣使送还斛斯政后隋军撤回”，发生在隋与高句丽的政治语境中。账本将其原因概括为连续远征未获成果，隋廷须应付国内起事和军队疲敝，高句丽以有限外交动作换取时间，并提示辽东战争暂止却未解决竞争，隋的动员信用和地方服从继续下滑。这类节点进入地方社会时，要经过官员、军队、书吏、商人、寺院、道路与家庭等不同中介；同一政策或战事在城市、边地和乡村不会具有相同的节奏。

核心来源为《隋书·炀帝纪下》；《隋书·东夷传·高丽》；《资治通鉴》隋纪（大业十年条）。它们可以较有把握地支持年序、制度公布、使节活动或特定地点的记载，却不能自动给出人口、财富、兵额、族属和内心动机的完整调查。账本的证据注记是“出征与撤退可靠；遣还人物是否构成实质和议及谈判条件不明”。因此，不能将官府的分类、战报的数字、宗教文本的愿望或后出的总括直接写成全体社会事实。

更合适的解释是，将这一事件看成资源、信息和权力重新连接的一环：命令是否抵达，取决于交通和地方协作；征发是否落实，取决于登记、市场和劳力；文化和宗教网络能否延续，也受战争、迁徙和赞助条件影响。现有材料只照亮其中一部分，叙事必须让区域差异与材料沉默继续可见。

条目\texttt{ST-615-YANMEN-SIEGE}所记的大业十一年（615）八月“始毕可汗围炀帝于雁门，隋廷以勤王与外交解围”，发生在隋与东突厥的政治语境中。账本将其原因概括为隋利用突厥内部分化而边防兵力被征辽牵制，始毕得以深入雁门，并提示皇帝受围暴露边防脆弱，勤王军与地方豪强的军事地位上升。这类节点进入地方社会时，要经过官员、军队、书吏、商人、寺院、道路与家庭等不同中介；同一政策或战事在城市、边地和乡村不会具有相同的节奏。

核心来源为《隋书·炀帝纪下》；《隋书·突厥传》；《北史·突厥传》。它们可以较有把握地支持年序、制度公布、使节活动或特定地点的记载，却不能自动给出人口、财富、兵额、族属和内心动机的完整调查。账本的证据注记是“围城事实可靠；义成公主传讯和解围细节多由汉文叙事单线保存”。因此，不能将官府的分类、战报的数字、宗教文本的愿望或后出的总括直接写成全体社会事实。

更合适的解释是，将这一事件看成资源、信息和权力重新连接的一环：命令是否抵达，取决于交通和地方协作；征发是否落实，取决于登记、市场和劳力；文化和宗教网络能否延续，也受战争、迁徙和赞助条件影响。现有材料只照亮其中一部分，叙事必须让区域差异与材料沉默继续可见。

条目\texttt{ST-617-LI-YUAN-TAIYUAN}所记的义宁元年（617）五月至十一月“李渊在太原起兵，入关后拥立代王杨侑，受封唐王”，发生在唐国与隋的政治语境中。账本将其原因概括为隋廷失去对北方军镇和财政节点的控制，李渊凭太原兵马、门阀及突厥外交取得资源，并提示关中出现能继承中央机构的竞争集团，长安成为唐建国基础。这类节点进入地方社会时，要经过官员、军队、书吏、商人、寺院、道路与家庭等不同中介；同一政策或战事在城市、边地和乡村不会具有相同的节奏。

核心来源为《旧唐书·高祖本纪》；《新唐书·高祖本纪》；《资治通鉴》隋纪（义宁元年条）。它们可以较有把握地支持年序、制度公布、使节活动或特定地点的记载，却不能自动给出人口、财富、兵额、族属和内心动机的完整调查。账本的证据注记是“起兵与入关序列可靠；与突厥结盟条款及李世民作用在唐代叙事中被重塑”。因此，不能将官府的分类、战报的数字、宗教文本的愿望或后出的总括直接写成全体社会事实。

更合适的解释是，将这一事件看成资源、信息和权力重新连接的一环：命令是否抵达，取决于交通和地方协作；征发是否落实，取决于登记、市场和劳力；文化和宗教网络能否延续，也受战争、迁徙和赞助条件影响。现有材料只照亮其中一部分，叙事必须让区域差异与材料沉默继续可见。

条目\texttt{ST-618-JIANGDU-COUP}所记的大业十四年（618）三月“宇文化及等发动江都兵变，杀隋炀帝并挟持隋廷北返”，发生在隋与骁果军的政治语境中。账本将其原因概括为骁果军滞留江都、归北无期且供给压力加重，宫廷对军队失去控制，并提示隋失去名义最高统治者，各地竞争者加速以新国号争夺合法性。这类节点进入地方社会时，要经过官员、军队、书吏、商人、寺院、道路与家庭等不同中介；同一政策或战事在城市、边地和乡村不会具有相同的节奏。

核心来源为《隋书·炀帝纪下》；《隋书·宇文化及传》；《资治通鉴》隋纪（大业十四年三月条）。它们可以较有把握地支持年序、制度公布、使节活动或特定地点的记载，却不能自动给出人口、财富、兵额、族属和内心动机的完整调查。账本的证据注记是“政变与炀帝被杀可靠；参与者责任、遗诏和宫廷细节多为唐初道德化叙述”。因此，不能将官府的分类、战报的数字、宗教文本的愿望或后出的总括直接写成全体社会事实。

更合适的解释是，将这一事件看成资源、信息和权力重新连接的一环：命令是否抵达，取决于交通和地方协作；征发是否落实，取决于登记、市场和劳力；文化和宗教网络能否延续，也受战争、迁徙和赞助条件影响。现有材料只照亮其中一部分，叙事必须让区域差异与材料沉默继续可见。

条目\texttt{ST-618-TANG-FOUNDING}所记的武德元年（618）五月“李渊受杨侑禅，建唐，改元武德”，发生在唐的政治语境中。账本将其原因概括为李渊已控制长安和关中官僚财政，隋炀帝死后代王政权更无独立军政基础，并提示唐取得中央名义与隋制遗产，但须在各地击败竞争政权。这类节点进入地方社会时，要经过官员、军队、书吏、商人、寺院、道路与家庭等不同中介；同一政策或战事在城市、边地和乡村不会具有相同的节奏。

核心来源为《旧唐书·高祖本纪》；《新唐书·高祖本纪》；《资治通鉴》唐纪（武德元年五月条）。它们可以较有把握地支持年序、制度公布、使节活动或特定地点的记载，却不能自动给出人口、财富、兵额、族属和内心动机的完整调查。账本的证据注记是“受禅日期可靠；唐修史对杨侑与隋末合法性的呈现服务新朝继统”。因此，不能将官府的分类、战报的数字、宗教文本的愿望或后出的总括直接写成全体社会事实。

更合适的解释是，将这一事件看成资源、信息和权力重新连接的一环：命令是否抵达，取决于交通和地方协作；征发是否落实，取决于登记、市场和劳力；文化和宗教网络能否延续，也受战争、迁徙和赞助条件影响。现有材料只照亮其中一部分，叙事必须让区域差异与材料沉默继续可见。

条目\texttt{ST-619-ZHENG-FOUNDING}所记的武德二年（619）“王世充废杨侗，自立国号郑，以洛阳为中心抗唐”，发生在郑与唐的政治语境中。账本将其原因概括为杨侗政权倚赖王世充军力，唐关中政权东进迫使洛阳形成独立资源中心，并提示中原出现郑、夏、唐竞争，洛阳和运河粮道成为关键目标。这类节点进入地方社会时，要经过官员、军队、书吏、商人、寺院、道路与家庭等不同中介；同一政策或战事在城市、边地和乡村不会具有相同的节奏。

核心来源为《旧唐书·王世充传》；《新唐书·王世充传》；《资治通鉴》唐纪（武德二年条）。它们可以较有把握地支持年序、制度公布、使节活动或特定地点的记载，却不能自动给出人口、财富、兵额、族属和内心动机的完整调查。账本的证据注记是“称帝和年次可靠；洛阳民众支持度主要被敌对史书以成败倒叙”。因此，不能将官府的分类、战报的数字、宗教文本的愿望或后出的总括直接写成全体社会事实。

更合适的解释是，将这一事件看成资源、信息和权力重新连接的一环：命令是否抵达，取决于交通和地方协作；征发是否落实，取决于登记、市场和劳力；文化和宗教网络能否延续，也受战争、迁徙和赞助条件影响。现有材料只照亮其中一部分，叙事必须让区域差异与材料沉默继续可见。

条目\texttt{ST-621-HULAO}所记的武德四年（621）五月“李世民在虎牢击败窦建德，王世充随后降唐，唐取洛阳”，发生在唐、郑与夏的政治语境中。账本将其原因概括为唐围洛阳迫使窦建德南援，唐军在虎牢阻止夏军与郑军会合，并提示唐夺取东都和中原仓储，华北竞争格局转向唐军可逐步收束。这类节点进入地方社会时，要经过官员、军队、书吏、商人、寺院、道路与家庭等不同中介；同一政策或战事在城市、边地和乡村不会具有相同的节奏。

核心来源为《旧唐书·太宗本纪上》；《旧唐书·窦建德传》；《新唐书·王世充传》；《资治通鉴》唐纪（武德四年条）。它们可以较有把握地支持年序、制度公布、使节活动或特定地点的记载，却不能自动给出人口、财富、兵额、族属和内心动机的完整调查。账本的证据注记是“战役和郑降可靠；兵力、突阵与战果数字多为唐方军功叙事”。因此，不能将官府的分类、战报的数字、宗教文本的愿望或后出的总括直接写成全体社会事实。

更合适的解释是，将这一事件看成资源、信息和权力重新连接的一环：命令是否抵达，取决于交通和地方协作；征发是否落实，取决于登记、市场和劳力；文化和宗教网络能否延续，也受战争、迁徙和赞助条件影响。现有材料只照亮其中一部分，叙事必须让区域差异与材料沉默继续可见。

条目\texttt{ST-622-LIU-HEITA}所记的武德五年（622）“刘黑闼在河北再起，败后被部将擒杀，唐重新控制河北主要州县”，发生在唐与河北反唐军的政治语境中。账本将其原因概括为唐处置夏地人事和征发引发反弹，窦建德旧部仍有地方网络与动员能力，并提示唐须以安抚、任官和驻军巩固河北，统一战争转入地方接管。这类节点进入地方社会时，要经过官员、军队、书吏、商人、寺院、道路与家庭等不同中介；同一政策或战事在城市、边地和乡村不会具有相同的节奏。

核心来源为《旧唐书·刘黑闼传》；《新唐书·刘黑闼传》；《资治通鉴》唐纪（武德五年条）。它们可以较有把握地支持年序、制度公布、使节活动或特定地点的记载，却不能自动给出人口、财富、兵额、族属和内心动机的完整调查。账本的证据注记是“刘黑闼败死可靠；地方倒向、军队构成和杀伤规模在史书中不完整”。因此，不能将官府的分类、战报的数字、宗教文本的愿望或后出的总括直接写成全体社会事实。

更合适的解释是，将这一事件看成资源、信息和权力重新连接的一环：命令是否抵达，取决于交通和地方协作；征发是否落实，取决于登记、市场和劳力；文化和宗教网络能否延续，也受战争、迁徙和赞助条件影响。现有材料只照亮其中一部分，叙事必须让区域差异与材料沉默继续可见。

条目\texttt{ST-626-XUANWU}所记的武德九年（626）六月“李世民在玄武门杀太子建成、齐王元吉，掌握禁军与朝政”，发生在唐的政治语境中。账本将其原因概括为东宫与秦王府长期争夺军功、人事和继承，皇帝未能建立可接受的裁决机制，并提示李世民控制继承与中枢，皇室暴力成为唐初制度阴影。这类节点进入地方社会时，要经过官员、军队、书吏、商人、寺院、道路与家庭等不同中介；同一政策或战事在城市、边地和乡村不会具有相同的节奏。

核心来源为《旧唐书·太宗本纪上》；《旧唐书·隐太子建成传》；《资治通鉴》唐纪（武德九年六月条）。它们可以较有把握地支持年序、制度公布、使节活动或特定地点的记载，却不能自动给出人口、财富、兵额、族属和内心动机的完整调查。账本的证据注记是“政变日期可靠；预谋、伏兵和言辞受太宗朝史料重塑，须与继承结果分开”。因此，不能将官府的分类、战报的数字、宗教文本的愿望或后出的总括直接写成全体社会事实。

更合适的解释是，将这一事件看成资源、信息和权力重新连接的一环：命令是否抵达，取决于交通和地方协作；征发是否落实，取决于登记、市场和劳力；文化和宗教网络能否延续，也受战争、迁徙和赞助条件影响。现有材料只照亮其中一部分，叙事必须让区域差异与材料沉默继续可见。

\subsection{文化、宗教与社群}

条目\texttt{ST-626-WEISHUI}所记的武德九年（626）八月“颉利可汗南至渭水北岸，李世民与其结盟退兵”，发生在唐与东突厥的政治语境中。账本将其原因概括为玄武门政变后唐军政调整未定，东突厥利用关中防务空隙南下，并提示唐以短期妥协避免京畿决战，随后整军、粮储和分化突厥成为优先事项。这类节点进入地方社会时，要经过官员、军队、书吏、商人、寺院、道路与家庭等不同中介；同一政策或战事在城市、边地和乡村不会具有相同的节奏。

核心来源为《旧唐书·太宗本纪上》；《旧唐书·突厥传上》；《资治通鉴》唐纪（武德九年八月条）。它们可以较有把握地支持年序、制度公布、使节活动或特定地点的记载，却不能自动给出人口、财富、兵额、族属和内心动机的完整调查。账本的证据注记是“会盟与退兵可靠；桥上谈判、财物数额及威慑效果主要见唐方叙事”。因此，不能将官府的分类、战报的数字、宗教文本的愿望或后出的总括直接写成全体社会事实。

更合适的解释是，将这一事件看成资源、信息和权力重新连接的一环：命令是否抵达，取决于交通和地方协作；征发是否落实，取决于登记、市场和劳力；文化和宗教网络能否延续，也受战争、迁徙和赞助条件影响。现有材料只照亮其中一部分，叙事必须让区域差异与材料沉默继续可见。

条目\texttt{ST-627-ZHENGUAN-CODE}所记的贞观元年（627）“唐廷修定贞观律令，形成初唐法制的规范文本”，发生在唐的政治语境中。账本将其原因概括为唐需在隋制基础上稳定刑名、官司和户婚秩序，并回应隋末苛法的政治记忆，并提示成文法提供共同术语，后续永徽律疏议将其进一步解释和定型。这类节点进入地方社会时，要经过官员、军队、书吏、商人、寺院、道路与家庭等不同中介；同一政策或战事在城市、边地和乡村不会具有相同的节奏。

核心来源为《旧唐书·刑法志》；《唐律疏议·名例》及疏议序；《通典·刑法》。它们可以较有把握地支持年序、制度公布、使节活动或特定地点的记载，却不能自动给出人口、财富、兵额、族属和内心动机的完整调查。账本的证据注记是“修定事实可靠；现存条文经后世保存，不能直接证明各州审判一致”。因此，不能将官府的分类、战报的数字、宗教文本的愿望或后出的总括直接写成全体社会事实。

更合适的解释是，将这一事件看成资源、信息和权力重新连接的一环：命令是否抵达，取决于交通和地方协作；征发是否落实，取决于登记、市场和劳力；文化和宗教网络能否延续，也受战争、迁徙和赞助条件影响。现有材料只照亮其中一部分，叙事必须让区域差异与材料沉默继续可见。

条目\texttt{ST-628-LIANG-SHIDU}所记的贞观二年（628）“梁师都部将杀梁师都降唐，朔方割据政权结束”，发生在唐与梁的政治语境中。账本将其原因概括为梁师都长期依东突厥支援，唐军北进与突厥内压削弱其外援和内部凝聚力，并提示唐完成关中北缘和河套门户控制，为反攻东突厥减少后方威胁。这类节点进入地方社会时，要经过官员、军队、书吏、商人、寺院、道路与家庭等不同中介；同一政策或战事在城市、边地和乡村不会具有相同的节奏。

核心来源为《旧唐书·梁师都传》；《新唐书·梁师都传》；《资治通鉴》唐纪（贞观二年条）。它们可以较有把握地支持年序、制度公布、使节活动或特定地点的记载，却不能自动给出人口、财富、兵额、族属和内心动机的完整调查。账本的证据注记是“降唐和年次可靠；部将倒戈的具体利益安排资料不足”。因此，不能将官府的分类、战报的数字、宗教文本的愿望或后出的总括直接写成全体社会事实。

更合适的解释是，将这一事件看成资源、信息和权力重新连接的一环：命令是否抵达，取决于交通和地方协作；征发是否落实，取决于登记、市场和劳力；文化和宗教网络能否延续，也受战争、迁徙和赞助条件影响。现有材料只照亮其中一部分，叙事必须让区域差异与材料沉默继续可见。

条目\texttt{ST-630-EASTERN-TURKS}所记的贞观四年（630）“李靖等破东突厥，颉利可汗被俘，东突厥汗国瓦解”，发生在唐与东突厥的政治语境中。账本将其原因概括为东突厥遭雪灾、内争与唐朝离间，唐军利用朔方和定襄方向实施冬季突袭，并提示唐取得北方纵深并接纳突厥部众，但草原政治未永久服从唐朝。这类节点进入地方社会时，要经过官员、军队、书吏、商人、寺院、道路与家庭等不同中介；同一政策或战事在城市、边地和乡村不会具有相同的节奏。

核心来源为《旧唐书·太宗本纪上》；《旧唐书·突厥传上》；《新唐书·突厥传上》；《资治通鉴》唐纪（贞观四年条）。它们可以较有把握地支持年序、制度公布、使节活动或特定地点的记载，却不能自动给出人口、财富、兵额、族属和内心动机的完整调查。账本的证据注记是“覆灭序列可靠；军队人数、俘获数与天可汗含义须区分唐礼仪和草原政治实际”。因此，不能将官府的分类、战报的数字、宗教文本的愿望或后出的总括直接写成全体社会事实。

更合适的解释是，将这一事件看成资源、信息和权力重新连接的一环：命令是否抵达，取决于交通和地方协作；征发是否落实，取决于登记、市场和劳力；文化和宗教网络能否延续，也受战争、迁徙和赞助条件影响。现有材料只照亮其中一部分，叙事必须让区域差异与材料沉默继续可见。

条目\texttt{ST-631-TURK-SETTLEMENT}所记的贞观五年（631）前后“唐将降附突厥部众安置于北方州府并设置羁縻管理，朝廷为此发生争议”，发生在唐与降附突厥诸部的政治语境中。账本将其原因概括为颉利败亡后部众需粮食、牧地和政治归属，唐廷在近塞安置与塞外扶立间权衡，并提示边疆治理由胜利转为供给、任官和身份分类问题，后突厥复兴显示其不稳。这类节点进入地方社会时，要经过官员、军队、书吏、商人、寺院、道路与家庭等不同中介；同一政策或战事在城市、边地和乡村不会具有相同的节奏。

核心来源为《旧唐书·突厥传上》；《新唐书·突厥传上》；《资治通鉴》唐纪（贞观五年条）。它们可以较有把握地支持年序、制度公布、使节活动或特定地点的记载，却不能自动给出人口、财富、兵额、族属和内心动机的完整调查。账本的证据注记是“政策大体可靠；安置人口、州界和内地化效果缺少逐户材料”。因此，不能将官府的分类、战报的数字、宗教文本的愿望或后出的总括直接写成全体社会事实。

更合适的解释是，将这一事件看成资源、信息和权力重新连接的一环：命令是否抵达，取决于交通和地方协作；征发是否落实，取决于登记、市场和劳力；文化和宗教网络能否延续，也受战争、迁徙和赞助条件影响。现有材料只照亮其中一部分，叙事必须让区域差异与材料沉默继续可见。

条目\texttt{ST-634-TUYUHUN}所记的贞观八年（634）“李靖等击吐谷浑，慕容伏允自杀，其子诺曷钵受唐册立”，发生在唐与吐谷浑的政治语境中。账本将其原因概括为吐谷浑与唐在河西商路、牧地和边民问题上冲突，唐担忧其与突厥联动，并提示唐在青海建立册封关系，却将吐谷浑置于吐蕃崛起与唐边防之间。这类节点进入地方社会时，要经过官员、军队、书吏、商人、寺院、道路与家庭等不同中介；同一政策或战事在城市、边地和乡村不会具有相同的节奏。

核心来源为《旧唐书·吐谷浑传》；《新唐书·吐谷浑传》；《资治通鉴》唐纪（贞观八年条）。它们可以较有把握地支持年序、制度公布、使节活动或特定地点的记载，却不能自动给出人口、财富、兵额、族属和内心动机的完整调查。账本的证据注记是“战役与册立可靠；追击路线、伤亡和吐谷浑内部选择主要由唐方记录”。因此，不能将官府的分类、战报的数字、宗教文本的愿望或后出的总括直接写成全体社会事实。

更合适的解释是，将这一事件看成资源、信息和权力重新连接的一环：命令是否抵达，取决于交通和地方协作；征发是否落实，取决于登记、市场和劳力；文化和宗教网络能否延续，也受战争、迁徙和赞助条件影响。现有材料只照亮其中一部分，叙事必须让区域差异与材料沉默继续可见。

条目\texttt{ST-635-NESTORIAN}所记的贞观九年（635）“阿罗本至长安，唐廷允许译经传教，景教进入都城宗教空间”，发生在唐与大秦景教社群的政治语境中。账本将其原因概括为丝路商旅和宗教人员进入长安，唐廷以审阅、翻译后的准入方式面对外来仪式，并提示景教获得有限制度承认，是观察都市宗教网络的窗口而非全国信仰比例。这类节点进入地方社会时，要经过官员、军队、书吏、商人、寺院、道路与家庭等不同中介；同一政策或战事在城市、边地和乡村不会具有相同的节奏。

核心来源为《唐会要·景教》；《大秦景教流行中国碑》（建中二年，781）；《册府元龟·外臣部》。它们可以较有把握地支持年序、制度公布、使节活动或特定地点的记载，却不能自动给出人口、财富、兵额、族属和内心动机的完整调查。账本的证据注记是“来华年代主要据781年碑刻及后出会要，缺少同时代官方原件”。因此，不能将官府的分类、战报的数字、宗教文本的愿望或后出的总括直接写成全体社会事实。

更合适的解释是，将这一事件看成资源、信息和权力重新连接的一环：命令是否抵达，取决于交通和地方协作；征发是否落实，取决于登记、市场和劳力；文化和宗教网络能否延续，也受战争、迁徙和赞助条件影响。现有材料只照亮其中一部分，叙事必须让区域差异与材料沉默继续可见。

条目\texttt{ST-638-SONGZHOU-TIBET}所记的贞观十二年（638）“松赞干布军至松州，唐军击退其攻势，双方转入婚盟谈判”，发生在唐与吐蕃的政治语境中。账本将其原因概括为吐蕃高原政权扩张并要求与唐建立婚姻关系，唐以边军胜利限定谈判条件，并提示唐蕃由接触转为制度化外交，婚盟未消除边疆资源竞争。这类节点进入地方社会时，要经过官员、军队、书吏、商人、寺院、道路与家庭等不同中介；同一政策或战事在城市、边地和乡村不会具有相同的节奏。

核心来源为《旧唐书·吐蕃传上》；《新唐书·吐蕃传上》；《资治通鉴》唐纪（贞观十二年条）。它们可以较有把握地支持年序、制度公布、使节活动或特定地点的记载，却不能自动给出人口、财富、兵额、族属和内心动机的完整调查。账本的证据注记是“战事与谈判转向可靠；吐蕃动机和兵力几乎只由唐方记录，须保留材料沉默”。因此，不能将官府的分类、战报的数字、宗教文本的愿望或后出的总括直接写成全体社会事实。

更合适的解释是，将这一事件看成资源、信息和权力重新连接的一环：命令是否抵达，取决于交通和地方协作；征发是否落实，取决于登记、市场和劳力；文化和宗教网络能否延续，也受战争、迁徙和赞助条件影响。现有材料只照亮其中一部分，叙事必须让区域差异与材料沉默继续可见。

条目\texttt{ST-640-GAOCHANG}所记的贞观十四年（640）“侯君集灭高昌，唐置西州并推进对天山南路的直接经营”，发生在唐与高昌的政治语境中。账本将其原因概括为高昌在唐与西突厥竞争中阻遏唐使，唐选择以驻军和州县取代间接关系，并提示西州成为东部天山支点，直接治理带来屯田、驻军和地方适应成本。这类节点进入地方社会时，要经过官员、军队、书吏、商人、寺院、道路与家庭等不同中介；同一政策或战事在城市、边地和乡村不会具有相同的节奏。

核心来源为《旧唐书·高昌传》；《新唐书·高昌传》；《资治通鉴》唐纪（贞观十四年条）；吐鲁番出土唐西州文书。它们可以较有把握地支持年序、制度公布、使节活动或特定地点的记载，却不能自动给出人口、财富、兵额、族属和内心动机的完整调查。账本的证据注记是“灭国和置州可靠；高昌地方反应及唐制落地须依吐鲁番多语文书和城址分期检验”。因此，不能将官府的分类、战报的数字、宗教文本的愿望或后出的总括直接写成全体社会事实。

更合适的解释是，将这一事件看成资源、信息和权力重新连接的一环：命令是否抵达，取决于交通和地方协作；征发是否落实，取决于登记、市场和劳力；文化和宗教网络能否延续，也受战争、迁徙和赞助条件影响。现有材料只照亮其中一部分，叙事必须让区域差异与材料沉默继续可见。

条目\texttt{ST-641-WENCHENG}所记的贞观十五年（641）“文成公主入吐蕃，与松赞干布成婚”，发生在唐与吐蕃的政治语境中。账本将其原因概括为松州冲突后双方需稳定使节、边贸和政治互认，唐以宗室女婚盟调节关系，并提示婚盟建立交往通道，但未改变高原与河西边界竞争的结构条件。这类节点进入地方社会时，要经过官员、军队、书吏、商人、寺院、道路与家庭等不同中介；同一政策或战事在城市、边地和乡村不会具有相同的节奏。

核心来源为《旧唐书·吐蕃传上》；《新唐书·吐蕃传上》；《资治通鉴》唐纪（贞观十五年条）。它们可以较有把握地支持年序、制度公布、使节活动或特定地点的记载，却不能自动给出人口、财富、兵额、族属和内心动机的完整调查。账本的证据注记是“婚姻事实可靠；关于带入技术、制度和宗教的宏大叙事多为后起建构”。因此，不能将官府的分类、战报的数字、宗教文本的愿望或后出的总括直接写成全体社会事实。

更合适的解释是，将这一事件看成资源、信息和权力重新连接的一环：命令是否抵达，取决于交通和地方协作；征发是否落实，取决于登记、市场和劳力；文化和宗教网络能否延续，也受战争、迁徙和赞助条件影响。现有材料只照亮其中一部分，叙事必须让区域差异与材料沉默继续可见。

条目\texttt{ST-643-CHENGQIAN}所记的贞观十七年（643）“太子李承乾谋反案发，被废为庶人，李治改立为太子”，发生在唐的政治语境中。账本将其原因概括为太子与魏王各自招纳官僚，太宗对继承人的干预加剧宫廷派系化，并提示李治成为继承人，高宗时期后族与皇后政治的进入条件改变。这类节点进入地方社会时，要经过官员、军队、书吏、商人、寺院、道路与家庭等不同中介；同一政策或战事在城市、边地和乡村不会具有相同的节奏。

核心来源为《旧唐书·太宗本纪下》；《旧唐书·承乾太子传》；《新唐书·太宗诸子传》。它们可以较有把握地支持年序、制度公布、使节活动或特定地点的记载，却不能自动给出人口、财富、兵额、族属和内心动机的完整调查。账本的证据注记是“废立事实可靠；谋反情节和魏王泰角色受宫廷记录与道德判断影响”。因此，不能将官府的分类、战报的数字、宗教文本的愿望或后出的总括直接写成全体社会事实。

更合适的解释是，将这一事件看成资源、信息和权力重新连接的一环：命令是否抵达，取决于交通和地方协作；征发是否落实，取决于登记、市场和劳力；文化和宗教网络能否延续，也受战争、迁徙和赞助条件影响。现有材料只照亮其中一部分，叙事必须让区域差异与材料沉默继续可见。

条目\texttt{ST-644-GOGURYEO-EXPEDITION}所记的贞观十八年（644）“唐以高句丽内乱和扶余璋争端为由，部署亲征辽东”，发生在唐与高句丽的政治语境中。账本将其原因概括为泉盖苏文政变及其与百济关系触动唐的区域秩序主张，太宗亦寻求东征威望，并提示东北亚战争重塑唐、新罗、百济和高句丽的联盟选择。这类节点进入地方社会时，要经过官员、军队、书吏、商人、寺院、道路与家庭等不同中介；同一政策或战事在城市、边地和乡村不会具有相同的节奏。

核心来源为《旧唐书·太宗本纪下》；《旧唐书·高丽传》；《资治通鉴》唐纪（贞观十八年条）。它们可以较有把握地支持年序、制度公布、使节活动或特定地点的记载，却不能自动给出人口、财富、兵额、族属和内心动机的完整调查。账本的证据注记是“出师决策可靠；唐方讨逆表述不能取代高句丽内部材料”。因此，不能将官府的分类、战报的数字、宗教文本的愿望或后出的总括直接写成全体社会事实。

更合适的解释是，将这一事件看成资源、信息和权力重新连接的一环：命令是否抵达，取决于交通和地方协作；征发是否落实，取决于登记、市场和劳力；文化和宗教网络能否延续，也受战争、迁徙和赞助条件影响。现有材料只照亮其中一部分，叙事必须让区域差异与材料沉默继续可见。

\subsection{环境、财政与社会后果}

条目\texttt{ST-645-LIAODONG-WITHDRAWAL}所记的贞观十九年（645）九月“唐军久攻安市城未克，因寒冷和补给压力撤出辽东”，发生在唐与高句丽的政治语境中。账本将其原因概括为唐军虽夺若干城堡，却无法在冬季前摧毁核心防御，高句丽依城防拖延补给，并提示亲征未能灭高句丽，唐改采与新罗协同、持续施压的长期战略。这类节点进入地方社会时，要经过官员、军队、书吏、商人、寺院、道路与家庭等不同中介；同一政策或战事在城市、边地和乡村不会具有相同的节奏。

核心来源为《旧唐书·太宗本纪下》；《旧唐书·高丽传》；《三国史记·高句丽本纪》。它们可以较有把握地支持年序、制度公布、使节活动或特定地点的记载，却不能自动给出人口、财富、兵额、族属和内心动机的完整调查。账本的证据注记是“撤军与安市城未下可靠；守城英雄和战损数字具有双方后世叙事层累”。因此，不能将官府的分类、战报的数字、宗教文本的愿望或后出的总括直接写成全体社会事实。

更合适的解释是，将这一事件看成资源、信息和权力重新连接的一环：命令是否抵达，取决于交通和地方协作；征发是否落实，取决于登记、市场和劳力；文化和宗教网络能否延续，也受战争、迁徙和赞助条件影响。现有材料只照亮其中一部分，叙事必须让区域差异与材料沉默继续可见。

条目\texttt{ST-646-XUEYANTUO}所记的贞观二十年（646）“唐联结回纥等部击溃薛延陀，漠北草原权力重组”，发生在唐与薛延陀的政治语境中。账本将其原因概括为东突厥覆亡后薛延陀扩张，唐转而支持竞争部落以防其控制漠北，并提示回纥等力量上升，唐的册封和军事干预成本增加。这类节点进入地方社会时，要经过官员、军队、书吏、商人、寺院、道路与家庭等不同中介；同一政策或战事在城市、边地和乡村不会具有相同的节奏。

核心来源为《旧唐书·薛延陀传》；《新唐书·薛延陀传》；《资治通鉴》唐纪（贞观二十年条）。它们可以较有把握地支持年序、制度公布、使节活动或特定地点的记载，却不能自动给出人口、财富、兵额、族属和内心动机的完整调查。账本的证据注记是“薛延陀败亡大体可靠；各部联盟的自称、人数和忠属关系主要经唐方转写”。因此，不能将官府的分类、战报的数字、宗教文本的愿望或后出的总括直接写成全体社会事实。

更合适的解释是，将这一事件看成资源、信息和权力重新连接的一环：命令是否抵达，取决于交通和地方协作；征发是否落实，取决于登记、市场和劳力；文化和宗教网络能否延续，也受战争、迁徙和赞助条件影响。现有材料只照亮其中一部分，叙事必须让区域差异与材料沉默继续可见。

条目\texttt{ST-648-KUCHA}所记的贞观二十二年（648）“阿史那社尔攻克龟兹，唐在西域经营安西四镇的基础扩大”，发生在唐与龟兹的政治语境中。账本将其原因概括为唐控制高昌后仍须突破龟兹等绿洲联盟和西突厥影响以保障通道，并提示军事胜利转为都护府镇戍体系，西域治理更依赖运输、翻译和地方中介。这类节点进入地方社会时，要经过官员、军队、书吏、商人、寺院、道路与家庭等不同中介；同一政策或战事在城市、边地和乡村不会具有相同的节奏。

核心来源为《旧唐书·龟兹传》；《旧唐书·阿史那社尔传》；《新唐书·西域传》；吐鲁番出土唐代军政文书。它们可以较有把握地支持年序、制度公布、使节活动或特定地点的记载，却不能自动给出人口、财富、兵额、族属和内心动机的完整调查。账本的证据注记是“战役可靠；四镇设置先后、城址与常驻兵额须以文书和遗址分层核验”。因此，不能将官府的分类、战报的数字、宗教文本的愿望或后出的总括直接写成全体社会事实。

更合适的解释是，将这一事件看成资源、信息和权力重新连接的一环：命令是否抵达，取决于交通和地方协作；征发是否落实，取决于登记、市场和劳力；文化和宗教网络能否延续，也受战争、迁徙和赞助条件影响。现有材料只照亮其中一部分，叙事必须让区域差异与材料沉默继续可见。

条目\texttt{ST-649-TAIZONG-DEATH}所记的贞观二十三年（649）五月“太宗卒，太子李治即位，是为高宗”，发生在唐的政治语境中。账本将其原因概括为李治已在承乾废立后确立储位，重臣、后族和边疆军政仍需重新协调，并提示高宗继承扩张中的帝国与官僚，后宫、外戚和宰相关系成为新焦点。这类节点进入地方社会时，要经过官员、军队、书吏、商人、寺院、道路与家庭等不同中介；同一政策或战事在城市、边地和乡村不会具有相同的节奏。

核心来源为《旧唐书·太宗本纪下》；《旧唐书·高宗本纪上》；《资治通鉴》唐纪（贞观二十三年条）。它们可以较有把握地支持年序、制度公布、使节活动或特定地点的记载，却不能自动给出人口、财富、兵额、族属和内心动机的完整调查。账本的证据注记是“死亡与即位日期可靠；晚年用兵、遗诏和继承安排的评价受后世史论影响”。因此，不能将官府的分类、战报的数字、宗教文本的愿望或后出的总括直接写成全体社会事实。

更合适的解释是，将这一事件看成资源、信息和权力重新连接的一环：命令是否抵达，取决于交通和地方协作；征发是否落实，取决于登记、市场和劳力；文化和宗教网络能否延续，也受战争、迁徙和赞助条件影响。现有材料只照亮其中一部分，叙事必须让区域差异与材料沉默继续可见。

条目\texttt{ST-653-TANGLU-SHUYI}所记的永徽四年（653）“长孙无忌等奉诏修成《唐律疏议》，为律文附加官方解释”，发生在唐的政治语境中。账本将其原因概括为贞观律令需随官制、礼法解释和案件经验统一，朝廷借重宰相集团完成注疏，并提示律与疏成为后世东亚法典参照，也凸显中央规范与地方实际的差异。这类节点进入地方社会时，要经过官员、军队、书吏、商人、寺院、道路与家庭等不同中介；同一政策或战事在城市、边地和乡村不会具有相同的节奏。

核心来源为《唐律疏议·序》；《旧唐书·刑法志》；《新唐书·刑法志》。它们可以较有把握地支持年序、制度公布、使节活动或特定地点的记载，却不能自动给出人口、财富、兵额、族属和内心动机的完整调查。账本的证据注记是“成书年代大体可靠；现存版本经传抄整理，律疏是规范解释而非审判统计”。因此，不能将官府的分类、战报的数字、宗教文本的愿望或后出的总括直接写成全体社会事实。

更合适的解释是，将这一事件看成资源、信息和权力重新连接的一环：命令是否抵达，取决于交通和地方协作；征发是否落实，取决于登记、市场和劳力；文化和宗教网络能否延续，也受战争、迁徙和赞助条件影响。现有材料只照亮其中一部分，叙事必须让区域差异与材料沉默继续可见。

条目\texttt{ST-655-WU-EMPRESS}所记的永徽六年（655）“高宗废王皇后、立武昭仪为皇后”，发生在唐的政治语境中。账本将其原因概括为高宗继位后后宫、舅族和宰相围绕继承及人事结盟，武昭仪获得皇帝支持，并提示皇后更替重组派系，武后由此进入影响诏令、人事和继承的权力位置。这类节点进入地方社会时，要经过官员、军队、书吏、商人、寺院、道路与家庭等不同中介；同一政策或战事在城市、边地和乡村不会具有相同的节奏。

核心来源为《旧唐书·高宗本纪上》；《旧唐书·则天皇后纪》；《新唐书·后妃传上》。它们可以较有把握地支持年序、制度公布、使节活动或特定地点的记载，却不能自动给出人口、财富、兵额、族属和内心动机的完整调查。账本的证据注记是“废立日期可靠；王皇后、萧淑妃与武后形象深受后代性别化叙事塑造”。因此，不能将官府的分类、战报的数字、宗教文本的愿望或后出的总括直接写成全体社会事实。

更合适的解释是，将这一事件看成资源、信息和权力重新连接的一环：命令是否抵达，取决于交通和地方协作；征发是否落实，取决于登记、市场和劳力；文化和宗教网络能否延续，也受战争、迁徙和赞助条件影响。现有材料只照亮其中一部分，叙事必须让区域差异与材料沉默继续可见。

条目\texttt{ST-657-WESTERN-TURKS}所记的显庆二年（657）“苏定方击阿史那贺鲁，唐推进都护与羁縻安排”，发生在唐与西突厥的政治语境中。账本将其原因概括为西突厥内争和唐在高昌、龟兹的驻军形成夹击，贺鲁与唐竞争绿洲中介，并提示唐西向影响达到高点，但长距离补给和地方联盟使都护体系不等于稳固领土。这类节点进入地方社会时，要经过官员、军队、书吏、商人、寺院、道路与家庭等不同中介；同一政策或战事在城市、边地和乡村不会具有相同的节奏。

核心来源为《旧唐书·苏定方传》；《旧唐书·突厥传下》；《新唐书·西突厥传》。它们可以较有把握地支持年序、制度公布、使节活动或特定地点的记载，却不能自动给出人口、财富、兵额、族属和内心动机的完整调查。账本的证据注记是“击败贺鲁可靠；对草原和中亚的实际控制范围、归附持续性不可由册封文辞推定”。因此，不能将官府的分类、战报的数字、宗教文本的愿望或后出的总括直接写成全体社会事实。

更合适的解释是，将这一事件看成资源、信息和权力重新连接的一环：命令是否抵达，取决于交通和地方协作；征发是否落实，取决于登记、市场和劳力；文化和宗教网络能否延续，也受战争、迁徙和赞助条件影响。现有材料只照亮其中一部分，叙事必须让区域差异与材料沉默继续可见。

条目\texttt{ST-660-BAEKJE}所记的显庆五年（660）“唐将苏定方与新罗军联合作战，百济王义慈降，百济覆亡”，发生在唐、新罗与百济的政治语境中。账本将其原因概括为高句丽战事受挫后唐借新罗开辟半岛战场，新罗则谋求打破百济高句丽联盟，并提示百济领土成为分配问题，唐置都督府的计划与新罗利益冲突加深。这类节点进入地方社会时，要经过官员、军队、书吏、商人、寺院、道路与家庭等不同中介；同一政策或战事在城市、边地和乡村不会具有相同的节奏。

核心来源为《旧唐书·百济传》；《三国史记·百济本纪》；《资治通鉴》唐纪（显庆五年条）。它们可以较有把握地支持年序、制度公布、使节活动或特定地点的记载，却不能自动给出人口、财富、兵额、族属和内心动机的完整调查。账本的证据注记是“百济覆亡可靠；唐、新罗两方对功劳、占领与降服的叙事立场不同”。因此，不能将官府的分类、战报的数字、宗教文本的愿望或后出的总括直接写成全体社会事实。

更合适的解释是，将这一事件看成资源、信息和权力重新连接的一环：命令是否抵达，取决于交通和地方协作；征发是否落实，取决于登记、市场和劳力；文化和宗教网络能否延续，也受战争、迁徙和赞助条件影响。现有材料只照亮其中一部分，叙事必须让区域差异与材料沉默继续可见。

条目\texttt{ST-663-BAEKGANG}所记的龙朔三年（663）八月“白江口海战中唐新罗联军击败倭国援助的百济复国军”，发生在唐、新罗、倭与百济复国军的政治语境中。账本将其原因概括为百济覆亡后遗民求倭援并据城复国，唐罗需切断海上补给，并提示倭与半岛军事关系调整，唐罗仍须面对高句丽与新罗的后续竞争。这类节点进入地方社会时，要经过官员、军队、书吏、商人、寺院、道路与家庭等不同中介；同一政策或战事在城市、边地和乡村不会具有相同的节奏。

核心来源为《旧唐书·刘仁轨传》；《三国史记·新罗本纪》；《日本书纪·天智二年》。它们可以较有把握地支持年序、制度公布、使节活动或特定地点的记载，却不能自动给出人口、财富、兵额、族属和内心动机的完整调查。账本的证据注记是“战役可靠；舰船、伤亡和倭军目的在中日韩文献中各有叙事差异”。因此，不能将官府的分类、战报的数字、宗教文本的愿望或后出的总括直接写成全体社会事实。

更合适的解释是，将这一事件看成资源、信息和权力重新连接的一环：命令是否抵达，取决于交通和地方协作；征发是否落实，取决于登记、市场和劳力；文化和宗教网络能否延续，也受战争、迁徙和赞助条件影响。现有材料只照亮其中一部分，叙事必须让区域差异与材料沉默继续可见。

条目\texttt{ST-664-XUANZANG-DEATH}所记的麟德元年（664）“玄奘卒于玉华宫，其译场和弟子网络继续传播佛典”，发生在唐与佛教僧团的政治语境中。账本将其原因概括为玄奘取经归来后获朝廷支持组织译场，僧侣、官员和写经网络维持译经事业，并提示个人声望转为寺院、文本和弟子网络，佛教知识生产继续嵌入都市空间。这类节点进入地方社会时，要经过官员、军队、书吏、商人、寺院、道路与家庭等不同中介；同一政策或战事在城市、边地和乡村不会具有相同的节奏。

核心来源为《大慈恩寺三藏法师传》；《续高僧传·玄奘传》；《旧唐书·方伎传》。它们可以较有把握地支持年序、制度公布、使节活动或特定地点的记载，却不能自动给出人口、财富、兵额、族属和内心动机的完整调查。账本的证据注记是“死亡年月大体可靠；生平神异和译经规模主要来自僧传，须同官方记录分读”。因此，不能将官府的分类、战报的数字、宗教文本的愿望或后出的总括直接写成全体社会事实。

更合适的解释是，将这一事件看成资源、信息和权力重新连接的一环：命令是否抵达，取决于交通和地方协作；征发是否落实，取决于登记、市场和劳力；文化和宗教网络能否延续，也受战争、迁徙和赞助条件影响。现有材料只照亮其中一部分，叙事必须让区域差异与材料沉默继续可见。

条目\texttt{ST-666-TAISHAN}所记的乾封元年（666）“高宗与武后封禅泰山，以国家礼仪宣示皇权与功业”，发生在唐的政治语境中。账本将其原因概括为唐在西突厥、百济等方向取胜，朝廷将边功、天命和皇后地位编入公共仪式，并提示封禅提高高宗武后共同出现的象征，成为武周合法性建构先例。这类节点进入地方社会时，要经过官员、军队、书吏、商人、寺院、道路与家庭等不同中介；同一政策或战事在城市、边地和乡村不会具有相同的节奏。

核心来源为《旧唐书·高宗本纪上》；《新唐书·高宗本纪》；《唐会要·封禅》。它们可以较有把握地支持年序、制度公布、使节活动或特定地点的记载，却不能自动给出人口、财富、兵额、族属和内心动机的完整调查。账本的证据注记是“封禅举行可靠；礼仪所称四海升平是政治表述，不能推断各地同等安定”。因此，不能将官府的分类、战报的数字、宗教文本的愿望或后出的总括直接写成全体社会事实。

更合适的解释是，将这一事件看成资源、信息和权力重新连接的一环：命令是否抵达，取决于交通和地方协作；征发是否落实，取决于登记、市场和劳力；文化和宗教网络能否延续，也受战争、迁徙和赞助条件影响。现有材料只照亮其中一部分，叙事必须让区域差异与材料沉默继续可见。

条目\texttt{ST-668-GOGURYEO-CONQUEST}所记的总章元年（668）“李勣与新罗军攻破平壤，高句丽覆亡，唐设置安东都护府”，发生在唐、新罗与高句丽的政治语境中。账本将其原因概括为高句丽受连年战争、内部权力变动及唐罗夹攻消耗，唐军以持续围攻取得突破，并提示唐试图在半岛北部设府，随即与新罗围绕旧地控制发生冲突。这类节点进入地方社会时，要经过官员、军队、书吏、商人、寺院、道路与家庭等不同中介；同一政策或战事在城市、边地和乡村不会具有相同的节奏。

核心来源为《旧唐书·高丽传》；《新唐书·高丽传》；《三国史记·新罗本纪》。它们可以较有把握地支持年序、制度公布、使节活动或特定地点的记载，却不能自动给出人口、财富、兵额、族属和内心动机的完整调查。账本的证据注记是“覆亡与置都护府可靠；人口迁徙、地方接收和有效辖区须同半岛材料互校”。因此，不能将官府的分类、战报的数字、宗教文本的愿望或后出的总括直接写成全体社会事实。

更合适的解释是，将这一事件看成资源、信息和权力重新连接的一环：命令是否抵达，取决于交通和地方协作；征发是否落实，取决于登记、市场和劳力；文化和宗教网络能否延续，也受战争、迁徙和赞助条件影响。现有材料只照亮其中一部分，叙事必须让区域差异与材料沉默继续可见。
